\section{Introduction}
\label{sec:introduction}

Let $(\Omega,\mathcal F,\mu)$ be a probability space and let $T$ be an
invertible measurable transformation with measurable inverse that preserves
$\mu$. The transformation is \emph{mixing} if
\[
 \mu(A\cap T^{-n}B)\longrightarrow\mu(A)\mu(B)
 \qquad(|n|\to\infty)
\]
for every $A,B\in\mathcal F$. For an integer $k\ge2$, it is
\emph{mixing of order $k$} if
\[
 \mu\left(\bigcap_{i=1}^k T^{-t_i}A_i\right)
 \longrightarrow\prod_{i=1}^k\mu(A_i)
\]
whenever $t_1<\cdots<t_k$ are integers and
$\min_{i<k}(t_{i+1}-t_i)\to\infty$, for every fixed choice of the
measurable sets $A_i$. Invariance allows us to take $t_1=0$.

\begin{theorem}
\label{thm:main}
Every invertible mixing probability-preserving transformation is mixing
of every finite order. More explicitly, under the assumptions above, for
every $k\ge3$ and every $A_1,\ldots,A_k\in\mathcal F$,
\[
 \mu\bigl(A_1\cap T^{-n_1}A_2\cap\cdots\cap
             T^{-(n_1+\cdots+n_{k-1})}A_k\bigr)
 \longrightarrow\prod_{i=1}^k\mu(A_i)
\]
as $\min(n_1,\ldots,n_{k-1})\to\infty$ through positive integers.
No standardness or countable-generation assumption is imposed on the
original probability space.
\end{theorem}

Corollary~\ref{cor:endomorphisms-flows} extends the conclusion to mixing
endomorphisms on arbitrary probability spaces and to mixing flows whose
Koopman operators are strongly continuous on $L^1$. These consequences
are stated and proved in Section~\ref{sec:consequences}.

Theorem~\ref{thm:main} answers Rokhlin's multiple-mixing problem
affirmatively for one transformation. Rokhlin~\cite{Rokhlin1949}
introduced higher-order mixing and proved it for ergodic surjective
continuous endomorphisms of compact commutative groups with Haar
probability. The question whether ordinary mixing alone implies the next
order is presented as Rokhlin's 1949 question in de la Rue's
account~\cite[Section~1.1]{deLaRue2006}; Ryzhikov surveys the all-orders
formulation in~\cite[Section~1]{Ryzhikov2024}. We count the order by the
number of sets, as de la Rue does: order $k$ here corresponds to
multiplicity $k-1$ in Rokhlin and Ryzhikov. Neither a group structure nor
any other special form of the transformation is assumed in our theorem.

The distinction between one time generator and several is essential.
Ledrappier's zero-entropy, mixing $\Z^2$ example is not threefold
mixing~\cite{Ledrappier1978}; de la Rue explains its three-dot relation
and why a direct one-dimensional analogue fails
\cite[Sections~1.2 and~3.4, preprint version]{deLaRue2006}. On the positive side,
Kalikow proved the twofold-to-threefold implication for rank-one
transformations~\cite{Kalikow1984}. Host obtained all orders under a
singular-spectrum hypothesis~\cite{Host1991}, and Ryzhikov obtained
all-orders conclusions for mixing finite-rank automorphisms
\cite{Ryzhikov1993}. More recently, Kanigowski and Ravotti proved
that mixing locally uniformly shearing flows are mixing of all orders
\cite[Theorem~6.8]{KanigowskiRavotti2024}.
These results require structure not present in a general mixing
transformation. For strongly mixing actions of countable abelian groups
on separable probability spaces,
Bergelson and Zelada proved that the higher-order good layouts form a
strong Ramsey-large set~\cite[Theorem~1.11]{BergelsonZelada2024}; this
does not assert convergence along every diverging-gap layout required
here.

The language of joinings, originating with Furstenberg's work on
disjointness~\cite{Furstenberg1967}, makes a failure of higher-order
mixing visible as a nonproduct limit law with independent proper
marginals. Pairwise independence alone does not force a joining to be
product. For an aperiodic finite-alphabet process, Janvresse and
de la Rue show that a pairwise-independent joining in which one
coordinate is a continuous function of the other two forces positive
entropy~\cite{JanvresseDeLaRue2008}.
For the twofold/threefold question, de la Rue records both the limit
joining and a zero-entropy reduction
\cite[Propositions~3.1 and~3.2, preprint version]{deLaRue2006}.
At arbitrary order, an antecedent Pinsker-factor reduction is described
by Ryzhikov~\cite[p.~251]{Ryzhikov1997}, who attributes it to a lemma of
Thouvenot~\cite{Thouvenot1975Pinsker}. We prove below the exact
finite-alphabet, least-failure form needed here, including the passage
from a possibly nonstandard original probability space. The new
obstacle after that reduction is to control all residual interactions
of the proper-independent limit law, not merely a selected joining.
We therefore use an ordered family of array laws and prove the later
operator and cancellation steps directly. Markov-operator methods for
joinings and intertwinings provide a related precedent
\cite{Ryzhikov1994,Ryzhikov1997}; the rectangle and exact-line
statements used here are established under our specific hypotheses.

\subsection*{The structural argument}

Suppose that $k$ is the least failing order, and put $I=\{1,\ldots,k\}$.
The reduction gives a zero-entropy finite-alphabet mixing process $X$ and
centered one-site functions $f_1,\ldots,f_k$ with product correlations
bounded below by a fixed positive constant along a sequence of diverging
layouts. A single limit joining retains this failure, but independence
of all its proper marginals does not force it to be product. For each
$n\ge1$, repeated limits at sums of $n$ layouts give a law $\rho_n$
on $X^{I^n}$, with the layout indices sent to infinity in a specified
order. The law $\rho_1$ retains the correlation witness; $\rho_2$ is
an ordered square whose proper families of last-axis planes are
independent. Higher levels allow one axis to become part of a local
state while preserving that conditional-product property. No symmetry
under permutations of array axes is assumed.

We enlarge all these laws together. A local state $y$ retains an
$X$-coordinate $x(y)$ and a parameter $u(y)$; proper last-axis planes
remain independent even after conditioning on the whole parameter array.
A conditional-expectation energy construction, using the method of
sated extensions \cite{Austin2015Sated}, makes the extension saturated:
richer data in any further admissible extension do not
improve the conditional expectation of an old test. This permits a first
axis to be included in the local state while selected states on that axis
are included in the parameter. The main conditioning argument proves that
these extensions still have the prescribed relative independence. In a
square array, saturation consequently gives exact locality of row and
column conditional laws.

The next step turns that locality into a description by orthogonal lines
of functions. At a coordinate of a joint law, consider the closed span of
the conditional expectations of functions of the other coordinates.
We call this its active space. Row and column conditional-expectation
operators in a square satisfy two rectangle identities. The identities
force fixed positive parts on the intersections of the active spaces;
their polar parts form an abelian unitary group. This yields lines spanned
by modulus-one functions that can be completed along either a row or a
column so that their product is constant. The converse inclusion requires
both incident operator identities and is proved as part of the same argument.

Here the sated-extension method is implemented in an array category
whose specified conditional product laws survive countable inverse
limits. It does not require mixing of the enlarged system. The operator step
extracts multiplicatively matched orthogonal lines from concrete kernel
identities. It uses dense-range cancellation of positive operators and
does not require their inverses to be bounded.

\subsection*{From lines to the contradiction}

The parameter geometry of the square makes the active space in each
individual fibre equal to the span of its exact lines. We enumerate these
lines measurably and regard multiplication modulo constant phases as a
countable abelian group. The row laws match these groups by isomorphisms.
A counting-measure lift of the parameter laws partitions each group into
finite blocks of constant cardinality. Time permutes the blocks, and on
an abstract copy of the group this permutation is implemented by a single
automorphism. Only the motion of whole blocks is used; a uniform rule for
the time evolution of individual labels inside a block is unnecessary.

Finally, a label orbit is either dissociated on arithmetic progressions
or polynomial on such progressions. In the first case an exponential
estimate for the label functions combines with the subexponential number
of likely words in a zero-entropy process. In the second case a Hilbert-space
averaging argument reduces polynomial degree, with ordinary mixing providing
the base case. These estimates allow an arbitrary coupling of the process
and the label functions. They show that the average absolute correlation
with each moving finite block tends to zero, contradicting a positive
quantity preserved by time.

The averages appear only in this contradiction inside a hypothetical
counterexample. The conclusion of Theorem~\ref{thm:main} remains convergence
along every diverging positive-gap layout. The proof uses one time
transformation throughout; the auxiliary array indices are not additional
time generators.

Section~\ref{sec:reduction} constructs the zero-entropy witness.
Section~\ref{sec:arrays} proves saturation and square locality.
Section~\ref{sec:operators} establishes the operator description of exact
lines. Section~\ref{sec:labels} obtains measurable labels and finite blocks.
Section~\ref{sec:cancellation} proves the one-dimensional cancellation
result, Section~\ref{sec:main-proof} assembles the contradiction, and
Section~\ref{sec:consequences} proves the endomorphism and flow consequences.

\subsection*{Conventions}

We write $\T=\{z\in\C:\abs z=1\}$ for the circle group under
multiplication. Hilbert spaces are complex, with inner products linear in
the second argument, and all spans of subspaces are closed spans. Conditional laws
and identities of factors are understood modulo null sets. After passing
to the finite-alphabet factor, all random-element spaces and subsequent
extensions are standard Borel spaces, with completion when needed.
Regular conditional laws can therefore be chosen using countable
determining classes. Assertions needed in finitely many orientations or
countably many time translates are imposed on a common conull set.
We use the standard conditional-expectation convergence theorems, the
finite-partition entropy chain rule, and the spectral theorem for compact
positive operators. Their applications below take place on the indicated
probability spaces or separable Hilbert spaces.
